\chapter{Skin Ulcer}

\section*{Definition}
A skin ulcer is a full-thickness defect of the epidermis and dermis that fails to progress through the normal healing sequence, persisting for more than 2 weeks. Chronic ulcers are classified by causes: venous, arterial, neuropathic (diabetic), or pressure-related.

\section*{What Happens in Your Body}
Chronic ulceration results from impaired tissue perfusion, persistent inflammation, and dysregulated wound healing. Venous ulcers arise from ambulatory venous hypertension causing fibrin cuff deposition and leukocyte trapping. Arterial ulcers reflect critical limb reduced blood flow from atherosclerotic occlusion. Neuropathic ulcers in diabetes involve sensory loss, motor deformity, and autonomic dysfunction leading to repetitive trauma and unperceived injury. Pressure ulcers develop from sustained mechanical loading causing reduced blood flow--reperfusion injury.

\section*{Causes}
\begin{itemize}
  \item \textbf{Venous:} Chronic venous insufficiency, deep vein thrombosis, varicose veins, obesity, immobility
  \item \textbf{Arterial:} Peripheral arterial disease, atherosclerosis, thromboangiitis obliterans, vasculitis
  \item \textbf{Neuropathic/diabetic:} Diabetes mellitus with peripheral neuropathy, Charcot foot
  \item \textbf{Pressure:} Immobility, spinal cord injury, prolonged bed rest, malnutrition, incontinence
  \item \textbf{Other:} Vasculitis (livedoid vasculopathy, pyoderma gangrenosum), infection (ecthyma, leishmaniasis), malignancy (basal cell carcinoma), sickle cell disease, calciphylaxis
\end{itemize}

\section*{When to See a Doctor}
See your doctor if:
\begin{itemize}
  \item New ulcer in you without clear causes or with atypical features
  \item Signs of infection: purulent drainage, increasing pain, surrounding cellulitis, fever
  \item Exposed bone, tendon, or joint (may indicate osteomyelitis)
  \item Rapid enlargement, necrotic black eschar, or malodorous discharge
  \item Failure to show signs of healing after 4--6 weeks of standard care
\end{itemize}

\section*{Related Symptoms}
\begin{itemize}
  \item Gangrene, osteomyelitis, cellulitis, lymphedema, peripheral neuropathy, claudication
\end{itemize}
\textbf{Important:} Wound assessment must include vascular evaluation (ABI, Doppler waveforms) before compression therapy, as applying compression to an arterial ulcer can worsen ischemia and cause limb loss.

\section*{Self-Care / Home Management}
\begin{itemize}
  \item Cleanse with normal saline or mild soap and water at each dressing change
  \item Apply appropriate moist wound dressing (hydrocolloid, foam, alginate, hydrogel) based on exudate level
  \item Offload pressure: special mattresses, wheelchair cushions, heel protectors, total contact casting for diabetic foot ulcers
\end{itemize}

\section*{How Your Doctor Will Evaluate You}
\begin{itemize}
  \item Determine causes: history (pain pattern, claudication, swelling, diabetes), location, wound bed appearance, and surrounding skin
  \item Vascular assessment: Ankle--brachial index (ABI), toe--brachial index (TBI), Doppler waveforms, venous duplex
  \item Wound culture (deep tissue or biopsy, not swab) if infection suspected; X-ray or MRI for osteomyelitis
  \item Biopsy of non-healing ulcer edge to rule out malignancy or specific inflammatory conditions
\end{itemize}

\section*{Treatment Options}
\begin{itemize}
  \item Venous: Compression therapy (multi-layer, 30--40 mmHg), leg elevation, pentoxifylline, wound removal of dead tissue
  \item Arterial: Revascularization (angioplasty, bypass), risk factor modification, antiplatelet therapy
  \item Diabetic/neuropathic: Offloading (total contact cast, CAM walker), sharp removal of dead tissue, infection control, glycemic optimization
  \item Pressure: Pressure redistribution, repositioning every 2 hours, nutritional support (protein, vitamin C, zinc)
  \item Advanced therapies: Negative pressure wound therapy, bioengineered skin substitutes, growth factors (PDGF), hyperbaric oxygen
\end{itemize}

\section*{References}
\begin{thebibliography}{99}
\bibitem{ref1} Singer AJ, Clark RAF. "Cutaneous wound healing." \textit{New England Journal of Medicine}, 1999;341(10):738--746.
\bibitem{ref2} Falanga V. "Wound healing and its impairment in the diabetic foot." \textit{The Lancet}, 2005;366(9498):1736--1743.
\bibitem{ref3} Sen CK, Gordillo GM, Roy S, et al. "Human skin wounds: A major and snowballing threat to public health and the economy." \textit{Wound Repair and Regeneration}, 2009;17(6):763--771.
\bibitem{ref4} Brem H, Tomic-Canic M, Tarnovskaya A, et al. "Healing of chronic wounds: A review of the literature." \textit{JAMA Dermatology}, 2007;143(11):1445--1449.
\bibitem{ref5} Moffatt CJ, Franks PJ, Doherty DC, et al. "Prevalence of leg ulceration in a London population." \textit{The Lancet}, 2004;363(9425):1501--1504.
\bibitem{ref6} O'Meara S, Cullum NA, Majid M, et al. "Systematic reviews of wound care management: Dressings and topical agents used in the healing of chronic wounds." \textit{Health Technology Assessment}, 2008;12(19):1--242.
\end{thebibliography}
